% ============================================
% DOC-17
% Gerado de molambudos.md — build_miolo.py
% ============================================
\section*{\textbf{DOC-17} --- Carta de um Soldado}
\markboth{DOC-17 --- Carta de um Soldado}{}

\noindent \textit{1915 --- o curral visto de dentro}

\noindent \textbf{Reconstituição ficcional --- documento narrativo}
\textbf{1915 --- Carta ficcional, sem equivalente arquivístico conhecido}

\vspace{0.5em}\noindent\rule{\textwidth}{0.5pt}\vspace{0.5em}

\noindent \textit{Carta ficcional atribuída ao soldado Antônio Alves Neto, lotado no Campo do Alagadiço, próximo a Fortaleza.}

\noindent \textit{Reconstituição ficcional endereçada a sua esposa, Maria, em Sobral; não foi localizada em arquivo militar.}

\vspace{0.5em}\noindent\rule{\textwidth}{0.5pt}\vspace{0.5em}

\noindent \textit{Maria,}

\noindent \textit{Estou escrevendo esta carta na tenda do curral, mas não sei se vou conseguir enviar. A situação aqui é pior do que o coronel conta para a imprensa.}

\noindent \textit{Todo dia chegam mais retirantes. Famílias inteiras. Gente que andou dias sem comida nem água. As crianças chegam com os olhos fundos e a barriga inchada. As mulheres, quando chegam, são separadas e levadas para \textquotedblleft{}realocação\textquotedblright{}. Eu sei o que isso significa. Todos sabem.}

\noindent \textit{Os homens e as crianças ficam. Ficam atrás do arame farpado, esperando a comida que nunca vem.}

\noindent \textit{O governo manda boi, mas o boi vai para a nossa cozinha, não para a deles. Eles recebem um caldo ralo, uma vez por dia. É o suficiente para não morrerem de uma vez. É o suficiente para morrerem aos poucos.}

\noindent \textit{O pior são as crianças. Tem um menino aqui, não deve ter mais de 10 anos, que chegou sozinho. A mãe foi levada no caminhão. O pai morreu na estrada. Ele não chora. Ele só olha. O olho dele é amarelo. O direito. Ele disse que \textquotedblleft{}um bicho\textquotedblright{} mudou. Ele não chora, Maria. Ele só olha. E olhar dele... me atravessa.}

\noindent \textit{Na semana passada, eu vi um negócio na vala.}

\noindent \textit{A vala dos corpos, atrás do curral. Eu estava de guarda noturna e ouvi um barulho. Mastigação. Fui ver.}

\noindent \textit{Tem uma coisa lá dentro. Mora no meio dos corpos. É grande. Muito grande. Tem olhos amarelos que brilham no escuro. Ela não me atacou. Ela só olhou. E continuou comendo.}

\noindent \textit{Contei para o sargento. Ele disse que é \textquotedblleft{}imaginação\textquotedblright{} e \textquotedblleft{}fome\textquotedblright{}. Disse que a fome faz a gente ver coisa.}

\noindent \textit{Talvez seja. Talvez não.}

\noindent \textit{Mas, esta noite, quando eu fechar os olhos, vou ver os olhos amarelos.}

\noindent \textit{Eu vi.}

\noindent \textit{E, Maria, eles me viram também.}

\noindent \textit{--- Antonio}

\vspace{0.5em}
\noindent\textit{\textbf{↪ Rotas:} \rota{DOC-18} \rotasep \rota{MEM-06}}
